\documentclass[../Odyssey_Project.tex]{subfiles}

\begin{document}
\setcounter{section}{14}
\section{The Character Table}

\subsubsection*{Definition of the Character Table}
\begin{itemize}
    \item Any character of $G$ is a $\Z$-linear combination of the $r$ irreducible characters $\chi_1, \ldots, \chi_r$ (Theorem~\ref{thm:14.7}), where $r$ is the number of conjugacy classes of $G$ (Theorem~\ref{thm:14.3}).
    \item Each irreducible character is determined by its values on the $r$ conjugacy classes, so the irreducible characters are completely determined by an $r \times r$ array.
    \item The \emph{character table} of $G$ is this $r \times r$ array $\mathcal{X} = (\chi_i(g_j))_{1 \leq i,j \leq r}$, where $g_1, \ldots, g_r$ are conjugacy class representatives.
    \item The character table is well-defined only up to reordering of rows and columns.
    \item By convention, $g_1 = 1$, so the first column consists of the degrees of $G$.
    \item We write $\mathcal{X}$ in the form
    \[
        \begin{array}{c|cccc}
            & 1 & k_2 & \cdots & k_r \\
            & 1 & g_2 & \cdots & g_r \\
            \hline
            \chi_1 & 1 & 1 & \cdots & 1 \\
            \chi_2 & f_2 & \chi_2(g_2) & \cdots & \chi_2(g_r) \\
            \vdots & \vdots & \vdots & \ddots & \vdots \\
            \chi_r & f_r & \chi_r(g_2) & \cdots & \chi_r(g_r)
        \end{array}
    \]
    where the $f_i$ are the degrees of $G$, and $k_i = |G : C_G(g_i)|$ is the order of the conjugacy class of $g_i$ (Proposition~\ref{prop:3.13}).
    \item By convention $\chi_1$ is always the principal character.
\end{itemize}

\subsubsection*{Orthogonality of Rows}

\begin{namedtheorem}[Row Orthogonality Theorem]
\label{thm:row-orthogonality}
$(\chi_i, \chi_j) = \delta_{ij}$ for any $i$ and $j$.
\end{namedtheorem}
\begin{proof}
    Let $S_1,\ldots,S_r$ be the distinct simple $\C G$-modules, ordered so that $S_i$ affords the character $\chi_i$. By Theorem~\ref{thm:14.12}, we have
    \[
        (\chi_i,\chi_j)=\dim_{\C}\Hom_{\C G}(S_i,S_j)
    \]
    for every $i$ and $j$. If $i=j$, then this Hom-space is $\End_{\C G}(S_i)$, which is isomorphic to $\C$ by Lemma~\ref{lem:13.14}; hence its dimension is $1$. If $i\neq j$, then $S_i$ and $S_j$ are non-isomorphic simple modules. By \nameref{lem:schur}, any non-zero map $S_i\to S_j$ would be an isomorphism, so no such non-zero map exists. Thus $\Hom_{\C G}(S_i,S_j)=0$ when $i\neq j$.\\
    Therefore $(\chi_i,\chi_j)$ is $1$ when $i=j$ and $0$ otherwise, which is exactly $\delta_{ij}$.
\end{proof}

In other words, this theorem asserts that
\[
    \delta_{ij} = \frac{1}{|G|} \sum_{g \in G} \chi_i(g)\overline{\chi_j(g)} = \frac{1}{|G|} \sum_{t=1}^{r} k_t \chi_i(g_t)\overline{\chi_j(g_t)}
\]
for any $i$ and $j$, where the $g_t$ are conjugacy class representatives and the $k_t$ are the orders of the conjugacy classes. The rows of the character table, considered as vectors in $\C^r$, are orthogonal with respect to an inner product that differs slightly from the standard one.

\begin{corollary}
\label{cor:15.1}
The irreducible characters of $G$ form an orthonormal basis for the vector space of class functions on $G$.
\end{corollary}
\begin{proof}
    Proposition~\ref{prop:14.10} gives that the irreducible characters form a basis for the vector space of class functions on $G$. By \nameref{thm:row-orthogonality}, their pairwise inner products satisfy $(\chi_i,\chi_j)=\delta_{ij}$. Thus this basis is orthonormal with respect to the inner product on class functions.
\end{proof}

\begin{corollary}
\label{cor:15.2}
If $\alpha = \sum_i a_i \chi_i$ and $\beta = \sum_j b_j \chi_j$ are virtual characters of $G$, then $(\alpha, \beta) = \sum_i a_i b_i$.
\end{corollary}
\begin{proof}
    Since $\alpha$ and $\beta$ are virtual characters, the coefficients $a_i$ and $b_j$ are integers. In particular, conjugating the coefficient $b_j$ does not change it. Using linearity in the first argument, conjugate-linearity in the second argument, and then row orthogonality, we get
    \[
        (\alpha,\beta)
        =
        \left(\sum_i a_i\chi_i,\sum_j b_j\chi_j\right)
        =
        \sum_i\sum_j a_i b_j(\chi_i,\chi_j)
        =
        \sum_i\sum_j a_i b_j\delta_{ij}
        =
        \sum_i a_i b_i.
    \]
\end{proof}

\begin{corollary}
\label{cor:15.3}
If $\alpha$ is a character of $G$ and $n \in \{1, 2, 3\}$, then $(\alpha, \alpha) = n$ iff $\alpha$ is a sum of $n$ distinct irreducible characters.
\end{corollary}
\begin{proof}
    Write
    \[
        \alpha=\sum_{i=1}^{r}a_i\chi_i
    \]
    as in Theorem~\ref{thm:14.7}. Since $\alpha$ is an ordinary character, the coefficients $a_i$ are non-negative integers. By Corollary~\ref{cor:15.2},
    \[
        (\alpha,\alpha)=\sum_{i=1}^{r}a_i^2.
    \]
    Suppose first that $(\alpha,\alpha)=n$ with $n\in\{1,2,3\}$. No coefficient $a_i$ can be at least $2$, because then $a_i^2\geq4>n$. Hence every $a_i$ is either $0$ or $1$, and exactly $n$ of them are equal to $1$. Thus $\alpha$ is the sum of exactly $n$ distinct irreducible characters.\\
    Conversely, if $\alpha$ is the sum of exactly $n$ distinct irreducible characters, then its coefficients in the irreducible-character basis are $1$ for those $n$ characters and $0$ for all the others. Corollary~\ref{cor:15.2} gives $(\alpha,\alpha)=n$.
\end{proof}

\begin{corollary}
\label{cor:15.4}
If $\alpha$ is a virtual character of $G$, then each $\chi_j$ appears with coefficient $(\alpha, \chi_j)$ in the unique expression of $\alpha$ as a $\Z$-linear combination of the irreducible characters of $G$.
\end{corollary}
\begin{proof}
    Write $\alpha=\sum_i a_i\chi_i$, with each $a_i\in\Z$. Applying Corollary~\ref{cor:15.2} to $\alpha$ and $\chi_j$ gives
    \[
        (\alpha,\chi_j)=a_j.
    \]
    Thus the inner product with $\chi_j$ reads off exactly the coefficient of $\chi_j$ in the irreducible-character expansion of $\alpha$.
\end{proof}

\begin{proposition}
\label{prop:15.5}
If $\alpha$ is a linear character of $G$ and $\chi$ is an irreducible character of $G$, then $\alpha\chi$ is an irreducible character of $G$.
\end{proposition}
\begin{proof}
    First, $\alpha\chi$ is a character by part (i) of Proposition~\ref{prop:14.8}, since it is the character of the tensor product of modules affording $\alpha$ and $\chi$. We only need to show that it is irreducible.\\
    Since $\alpha$ is linear, $\alpha(g)$ is a root of unity for every $g\in G$ by part (iii) of Proposition~\ref{prop:14.4}. Hence
    \[
        \alpha(g)\overline{\alpha(g)}=|\alpha(g)|^2=1
    \]
    for every $g\in G$. Therefore
    \begin{align*}
        (\alpha\chi,\alpha\chi)
        &=\frac{1}{|G|}\sum_{g\in G}\alpha(g)\chi(g)\overline{\alpha(g)\chi(g)} \\
        &=\frac{1}{|G|}\sum_{g\in G}\alpha(g)\overline{\alpha(g)}\chi(g)\overline{\chi(g)} \\
        &=\frac{1}{|G|}\sum_{g\in G}\chi(g)\overline{\chi(g)}
        =
        (\chi,\chi)
        =
        1,
    \end{align*}
    where the last equality is the Row Orthogonality Theorem. Since $\alpha\chi$ is a character with inner product $1$ against itself, Corollary~\ref{cor:15.3} shows that $\alpha\chi$ is irreducible.
\end{proof}

\subsubsection*{Orthogonality of Columns}

\begin{namedtheorem}[Column Orthogonality Theorem]
\label{thm:col-orthogonality}
If $g_1, \ldots, g_r$ are a set of conjugacy class representatives of $G$, and $k_1, \ldots, k_r$ are the orders of the conjugacy classes, then for any $1 \leq i, j \leq r$ we have
\[
    \sum_{t=1}^{r} \chi_t(g_i)\overline{\chi_t(g_j)} = \frac{|G|}{k_i}\delta_{ij} = |C_G(g_i)|\delta_{ij}.
\]
\end{namedtheorem}
\begin{proof}
    Let $\mathcal{X}=(\chi_i(g_j))_{1\leq i,j\leq r}$ be the character table, and let $K$ be the $r\times r$ diagonal matrix whose diagonal entries are $k_1,\ldots,k_r$. We first compute the $(i,j)$-entry of $\mathcal{X}K\overline{\mathcal{X}}^{\,T}$:
    \[
        (\mathcal{X}K\overline{\mathcal{X}}^{\,T})_{ij}
        =
        \sum_{t=1}^{r}\chi_i(g_t)k_t\overline{\chi_j(g_t)}
        =
        \sum_{g\in G}\chi_i(g)\overline{\chi_j(g)}.
    \]
    The last equality simply rewrites the sum over $G$ as a sum over conjugacy classes, since $\chi_i$ and $\chi_j$ are constant on conjugacy classes. By \nameref{thm:row-orthogonality}, this is $|G|\delta_{ij}$. Hence
    \[
        \mathcal{X}K\overline{\mathcal{X}}^{\,T}=|G|I_r.
    \]
    If square matrices $A$ and $B$ satisfy $AB=cI_r$ with $c\neq0$, then $A$ is invertible, $B=cA^{-1}$, and therefore $BA=cI_r$. Applying this to $A=\mathcal{X}$ and $B=K\overline{\mathcal{X}}^{\,T}$ gives
    \[
        K\overline{\mathcal{X}}^{\,T}\mathcal{X}=|G|I_r.
    \]
    Now take the $(i,j)$-entry of this last matrix equation. We obtain
    \[
        k_i\sum_{t=1}^{r}\overline{\chi_t(g_i)}\chi_t(g_j)=|G|\delta_{ij}.
    \]
    Taking complex conjugates gives
    \[
        \sum_{t=1}^{r}\chi_t(g_i)\overline{\chi_t(g_j)}
        =
        \frac{|G|}{k_i}\delta_{ij}.
    \]
    Finally, $k_i=|G:C_G(g_i)|$, so $|G|/k_i=|C_G(g_i)|$.
\end{proof}

\subsubsection*{Characters and Quotient Groups}

\begin{lemma}
\label{lem:15.6}
Let $N \trianglelefteq G$, and let $U$ be a $\C(G/N)$-module. Then $U$ admits a canonical $\C G$-module structure, with a subspace of $U$ being a $\C G$-submodule iff it is a $\C(G/N)$-submodule. If $\psi$ is the character of the $\C(G/N)$-module $U$, then the character of the $\C G$-module $U$ is $\psi \circ \eta$, where $\eta\colon G \to G/N$ is the natural map.
\end{lemma}
\begin{proof}
    Define the action of $G$ on $U$ by
    \[
        gu=(gN)u
    \]
    for $g\in G$ and $u\in U$, where the action on the right is the given action of $G/N$ on $U$. This is a $\C G$-module structure because the quotient action already satisfies
    \[
        ((gh)N)u=(gN)((hN)u),
        \qquad
        N u=u.
    \]
    A subspace $W\leq U$ is stable under this $G$-action exactly when it is stable under the action of every coset $gN\in G/N$. Thus the $\C G$-submodules of $U$ are precisely the $\C(G/N)$-submodules of $U$.\\
    For the character statement, fix $g\in G$. The linear transformation of $U$ induced by $g$ under this inflated $G$-action is exactly the same linear transformation as the one induced by the coset $\eta(g)=gN$ under the original $G/N$-action. Therefore its trace is $\psi(\eta(g))$, so the character as a $\C G$-module is $\psi\circ\eta$.
\end{proof}

We define $K_\chi = \{x \in G \mid \chi(x) = \chi(1)\}$ for any character $\chi$ of $G$; we have $K_\chi \trianglelefteq G$ by part (vi) of Proposition~\ref{prop:14.4}. We call $K_\chi$ the \emph{kernel} of $\chi$, as it is the kernel of the corresponding representation. We write $K_i$ instead of $K_{\chi_i}$.

\begin{proposition}
\label{prop:15.7}
The normal subgroups of $G$ are exactly the sets of the form $\bigcap_{i \in I} K_i$ for some $I \subseteq \{1, \ldots, r\}$.
\end{proposition}
\begin{proof}
    First let $N \trianglelefteq G$. Let $U=\C(G/N)$ be the regular $\C(G/N)$-module, and let $\psi$ be its character. We use Lemma~\ref{lem:15.6} to view $U$ as a $\C G$-module, and we write $\chi=\psi\circ\eta$ for its character as a $\C G$-module. Since $\psi$ is the regular character of $G/N$, we have
    \[
        \psi(gN)=\psi(N)
        \quad\Longleftrightarrow\quad
        gN=N.
    \]
    Therefore $\chi(g)=\chi(1)$ iff $g\in N$, so $K_\chi=N$.\\
    Now write $\chi=\sum_i a_i\chi_i$, where the $a_i$ are non-negative integers. We claim that
    \[
        K_\chi=\bigcap_{a_i>0}K_i.
    \]
    If $g\in K_i$ for every $i$ with $a_i>0$, then
    \[
        \chi(g)=\sum_i a_i\chi_i(g)=\sum_i a_i\chi_i(1)=\chi(1),
    \]
    so $g\in K_\chi$. Conversely, suppose $g\in K_\chi$. Then $\chi(g)=\chi(1)$, and hence
    \[
        \chi(1)=|\chi(g)|
        \leq
        \sum_i a_i|\chi_i(g)|
        \leq
        \sum_i a_i\chi_i(1)
        =
        \chi(1),
    \]
    where the second inequality is part (v) of Proposition~\ref{prop:14.4}. Thus equality holds throughout. For every $i$ with $a_i>0$, equality forces $|\chi_i(g)|=\chi_i(1)$, and equality in the triangle inequality forces the summands $a_i\chi_i(g)$ to point in the same direction as the positive real number $\chi(1)$. Hence $\chi_i(g)=\chi_i(1)$ for every $i$ with $a_i>0$, so $g\in\bigcap_{a_i>0}K_i$. This proves that $N=K_\chi$ has the required form.\\
    Conversely, each $K_i$ is normal in $G$ by part (vi) of Proposition~\ref{prop:14.4}. Therefore any finite intersection $\bigcap_{i\in I}K_i$ is normal in $G$ by repeated use of Proposition~\ref{prop:1.7}. If $I=\emptyset$, the intersection is $G$, which is also normal.
\end{proof}

\begin{corollary}
\label{cor:15.8}
$G$ is simple iff the only irreducible character $\chi_i$ for which $\chi_i(g) = \chi_i(1)$ for some $1 \neq g \in G$ is the principal character $\chi_1$.
\end{corollary}
\begin{proof}
    Suppose first that $G$ is simple. If $\chi_i(g)=\chi_i(1)$ for some $1\neq g\in G$, then $g\in K_i$. Thus $K_i$ is a non-trivial normal subgroup of $G$, so $K_i=G$. If $K_i=G$, then the corresponding irreducible representation is trivial on every element of $G$, and hence the simple module must be one-dimensional and trivial. Thus $\chi_i=\chi_1$.\\
    Conversely, suppose that $G$ is not simple. Then there is a non-trivial proper normal subgroup $N \trianglelefteq G$. Choose $1\neq g\in N$. By Proposition~\ref{prop:15.7}, we may write
    \[
        N=\bigcap_{i\in I}K_i
    \]
    for some $I\subseteq\{1,\ldots,r\}$. Since $N\neq G$, at least one non-principal $K_i$ must occur in this intersection; otherwise the intersection is just $G$, because $K_1=G$. Choose such an $i>1$. Then $g\in N\leq K_i$, so $\chi_i(g)=\chi_i(1)$ for a non-principal irreducible character. This is exactly the failure of the stated condition.
\end{proof}

\begin{corollary}
\label{cor:15.9}
The character table of $G$ can be used to determine whether or not $G$ is solvable.
\end{corollary}
\begin{proof}
    From the character table we can read off each $K_i$: it is the union of exactly those conjugacy classes on which $\chi_i$ takes the value $\chi_i(1)$. By Proposition~\ref{prop:15.7}, all normal subgroups of $G$ are obtained by intersecting these $K_i$. Thus the character table determines the list of normal subgroups of $G$.\\
    The table also records the sizes of the conjugacy classes, so once a normal subgroup has been identified as a union of conjugacy classes, its order can be computed by adding the corresponding class sizes. Hence the table determines the inclusion relations among normal subgroups and the orders of all quotients appearing in normal series. We can therefore check whether there is a normal series whose successive quotients have prime-power order, i.e. are $p$-groups. By Corollary~\ref{cor:11.7}, this is equivalent to solvability.
\end{proof}

We define $Z_\chi = \{x \in G \mid |\chi(x)| = \chi(1)\}$ for any character $\chi$ of $G$. Again we write $Z_i$ instead of $Z_{\chi_i}$.

\begin{lemma}
\label{lem:15.10}
$Z_\chi \leq G$ for any character $\chi$ of $G$, and if in addition $\chi$ is irreducible, then $Z_\chi/K_\chi = Z(G/K_\chi)$.
\end{lemma}
\begin{proof}
    Let $U$ be a $\C G$-module affording $\chi$, and let $\rho\colon G\to\mathrm{GL}(U)$ be the corresponding representation. First we show that $Z_\chi$ is a subgroup. Certainly $1\in Z_\chi$. If $g\in Z_\chi$, then $g^{-1}\in Z_\chi$ by part (iv) of Proposition~\ref{prop:14.4}. Now suppose $g,h\in Z_\chi$. By part (iii) of Proposition~\ref{prop:14.4}, $\chi(g)$ is a sum of $\chi(1)$ roots of unity. The condition $|\chi(g)|=\chi(1)$ means that this sum has the largest possible absolute value, so all eigenvalues of $\rho(g)$ are equal to a single root of unity; say $\rho(g)=\lambda(g)I$. Similarly, $\rho(h)=\lambda(h)I$. Hence
    \[
        \rho(gh)=\rho(g)\rho(h)=\lambda(g)\lambda(h)I,
    \]
    so all eigenvalues of $\rho(gh)$ are equal and $|\chi(gh)|=\chi(1)$. Thus $gh\in Z_\chi$, and $Z_\chi\leq G$.\\
    For any $g\in Z_\chi$, the matrix $\rho(g)$ is scalar, so it lies in the center of $\rho(G)$. Also $K_\chi\leq Z_\chi$, and since $K_\chi=\ker\rho$, we have $\rho(G)\cong G/K_\chi$. This gives
    \[
        Z_\chi/K_\chi\leq Z(G/K_\chi).
    \]
    Now assume that $\chi$ is irreducible, so $U$ is simple. Let $gK_\chi\in Z(G/K_\chi)$. Then $\rho(g)$ commutes with $\rho(x)$ for every $x\in G$, so the map $u\mapsto gu$ is a $\C G$-endomorphism of $U$. By Lemma~\ref{lem:13.14}, this endomorphism is multiplication by some scalar $\mu\in\C$. Since $g$ has finite order, $\mu$ is a root of unity. Therefore
    \[
        \chi(g)=\mu\chi(1),
    \]
    so $|\chi(g)|=\chi(1)$ and $g\in Z_\chi$. This proves the reverse inclusion, and hence $Z_\chi/K_\chi=Z(G/K_\chi)$.
\end{proof}

\begin{corollary}
\label{cor:15.11}
If $G$ is non-abelian and simple, then $Z_i = 1$ for all $i > 1$.
\end{corollary}
\begin{proof}
    Let $i>1$. Since $G$ is simple and $\chi_i$ is not the principal character, Corollary~\ref{cor:15.8} gives $K_i=1$. Applying Lemma~\ref{lem:15.10} to the irreducible character $\chi_i$ gives
    \[
        Z_i/K_i=Z(G/K_i).
    \]
    Since $K_i=1$, this says $Z_i=Z(G)$. But $G$ is non-abelian and simple, so $Z(G)$ cannot be all of $G$ and hence must be $1$. Therefore $Z_i=1$ for all $i>1$.
\end{proof}

\begin{proposition}
\label{prop:15.12}
$Z(G) = \bigcap_{i=1}^{r} Z_i$.
\end{proposition}
\begin{proof}
    First let $z\in Z(G)$. We show that $z\in Z_i$ for every $i$. Since $z$ is central in $G$, the coset $zK_i$ is central in $G/K_i$. By Lemma~\ref{lem:15.10}, applied to the irreducible character $\chi_i$, we have
    \[
        Z_i/K_i=Z(G/K_i).
    \]
    Hence $zK_i\in Z_i/K_i$. Since $K_i\leq Z_i$, this implies $z\in Z_i$. Therefore
    \[
        Z(G)\leq\bigcap_{i=1}^{r}Z_i.
    \]
    Conversely, suppose that $g\in\bigcap_{i=1}^{r}Z_i$. Fix any $x\in G$. For each $i$, Lemma~\ref{lem:15.10} gives
    \[
        gK_i\in Z_i/K_i=Z(G/K_i).
    \]
    Thus $gK_i$ commutes with $xK_i$ in $G/K_i$, so
    \[
        [g,x]\in K_i
    \]
    for every $i$. By Proposition~\ref{prop:15.7}, the intersection $K_1\cap\cdots\cap K_r$ is $1$: indeed, the normal subgroup $1$ is an intersection of some of the $K_i$, and intersecting all of them can only make this subgroup smaller. Hence $[g,x]=1$. Since this holds for every $x\in G$, we have $g\in Z(G)$.\\
    Therefore $Z(G)=\bigcap_{i=1}^{r}Z_i$.
\end{proof}

\begin{lemma}
\label{lem:15.13}
If $N \trianglelefteq G$, then the irreducible characters of $G/N$ can be determined from those of $G$.
\end{lemma}
\begin{proof}
    Let $\chi$ be an irreducible character of $G$ with $N\leq K_\chi$, and let $U$ be a simple $\C G$-module affording $\chi$. Since every element of $N$ acts as the identity on $U$, we can define an action of $G/N$ on $U$ by
    \[
        (gN)u=gu.
    \]
    This is well-defined: if $gN=hN$, then $h^{-1}g\in N$, so $h^{-1}g$ acts trivially on $U$, and hence $gu=hu$. A $\C(G/N)$-submodule of $U$ is the same thing as a $\C G$-submodule under this action. Since $U$ is simple as a $\C G$-module, it is also simple as a $\C(G/N)$-module. Thus $\chi$ gives an irreducible character of $G/N$, whose value on $gN$ is $\chi(g)$.\\
    Conversely, let $\psi$ be an irreducible character of $G/N$. By Lemma~\ref{lem:15.6}, the inflated character $\psi\circ\eta$ is the character of a $\C G$-module, and the submodule statement in that lemma shows that this inflated module is still simple. Hence $\psi\circ\eta$ is an irreducible character of $G$. Also, every element of $N$ maps to the identity coset, so $N\leq K_{\psi\circ\eta}$.\\
    Therefore the irreducible characters of $G/N$ are exactly the characters obtained from irreducible characters $\chi_i$ of $G$ satisfying $N\leq K_i$. The corresponding character of $G/N$ has value
    \[
        gN\mapsto\chi_i(g)
    \]
    on the coset $gN$. This determines the number of irreducible characters of $G/N$ and their values on elements of $G/N$ from the irreducible characters of $G$. What it does not automatically give is the full conjugacy-class structure or the class sizes in $G/N$, but the character values themselves are determined.
\end{proof}

\begin{corollary}
\label{cor:15.14}
The character table of $G$ can be used to determine whether or not $G$ is nilpotent.
\end{corollary}
\begin{proof}
    Let
    \[
        1=\mathcal{Z}_0\leq\mathcal{Z}_1\leq\mathcal{Z}_2\leq\cdots
    \]
    be the upper central series of $G$, so that
    \[
        \mathcal{Z}_m/\mathcal{Z}_{m-1}=Z(G/\mathcal{Z}_{m-1})
    \]
    for each $m\geq1$. We claim that the character table lets us determine each term $\mathcal{Z}_m$ recursively. For $m=1$, Proposition~\ref{prop:15.12} gives
    \[
        \mathcal{Z}_1=Z(G)=\bigcap_{i=1}^{r}Z_i,
    \]
    and each $Z_i$ can be read off from the character table by checking where $|\chi_i(g)|=\chi_i(1)$.\\
    Now suppose that $\mathcal{Z}_{m-1}$ has already been determined. Lemma~\ref{lem:15.13} tells us how to determine the irreducible characters of the quotient $G/\mathcal{Z}_{m-1}$ from the irreducible characters of $G$. Applying Proposition~\ref{prop:15.12} inside this quotient determines
    \[
        Z(G/\mathcal{Z}_{m-1})=\mathcal{Z}_m/\mathcal{Z}_{m-1}.
    \]
    Taking the full preimage in $G$ then determines $\mathcal{Z}_m$. Thus the whole upper central series is determined from the character table. By Exercise~\ref{ex:11.10}, the group $G$ is nilpotent iff $\mathcal{Z}_s=G$ for some $s$. Therefore the character table can be used to determine whether or not $G$ is nilpotent.
\end{proof}

Observe that in Corollary~\ref{cor:15.9} we assumed that we knew both the irreducible characters of $G$ and the orders of the conjugacy classes, whereas in Corollary~\ref{cor:15.14} we did not need to know the orders of the classes. As evidenced by Lemma~\ref{lem:15.13}, there may be circumstances in which we can construct a character table without knowing the orders of the conjugacy classes. There is a theorem of G.\ Higman (see [15, p.\ 136]) which asserts that if we know the irreducible characters of $G$, but not the orders of the conjugacy classes, then we can at least determine the sets of prime divisors of the orders of the conjugacy classes.

\subsubsection*{Examples: Constructing Character Tables}

The remainder of this section consists of a series of examples in which we develop methods of finding characters and use these methods to construct the character tables of various groups.

\begin{example}
\label{ex:15.eg.1}
Let $G = \langle g \rangle \cong \Z_n$ for some $n \in \N$, and let $\lambda$ be a primitive $n$th root of unity. For each $1 \leq i \leq n$, let $V_i$ be a one-dimensional $\C$-vector space, and let $g$ act on $V_i$ by multiplication by $\lambda^{i-1}$; since $G$ is cyclic, this definition completely determines a $\C G$-module structure on $V_i$. Each $V_i$, being one-dimensional, is a simple $\C G$-module. If $\chi_i$ is the character of $V_i$, then $\chi_i(g) = \lambda^{i-1}$, and hence $\chi_i(g^a) = \lambda^{a(i-1)}$ for any $a$. The characters $\chi_1, \ldots, \chi_n$ are $n$ distinct linear characters of $G$, and since $G$ can have at most $|G| = n$ irreducible characters, we see that the $\chi_i$ are precisely the irreducible characters of $G$. Thus, the character table of $G$ is
\[
    \begin{array}{c|ccccc}
        & 1 & 1 & 1 & \cdots & 1 \\
        & 1 & g & g^2 & \cdots & g^{n-1} \\
        \hline
        \chi_1 & 1 & 1 & 1 & \cdots & 1 \\
        \chi_2 & 1 & \lambda & \lambda^2 & \cdots & \lambda^{n-1} \\
        \chi_3 & 1 & \lambda^2 & \lambda^4 & \cdots & \lambda^{n-2} \\
        \vdots & \vdots & \vdots & \vdots & \ddots & \vdots \\
        \chi_n & 1 & \lambda^{n-1} & \lambda^{n-2} & \cdots & \lambda
    \end{array}
\]
Observe that the set $\{\chi_1, \ldots, \chi_n\}$ is a cyclic group under multiplication of characters, with generator $\chi_2$; we have $\chi_i = \chi_2^{i-1}$ for any $i$.
\end{example}

\begin{example}
\label{ex:15.eg.2}
Let $G$ and $H$ be groups with respective irreducible characters $\chi_1, \ldots, \chi_r$ and $\psi_1, \ldots, \psi_s$. We wish to determine the irreducible characters of $G \times H$. We see that two elements $(x, y)$ and $(x', y')$ of $G \times H$ are conjugate iff $x$ and $x'$ are conjugate in $G$ and $y$ and $y'$ are conjugate in $H$. Since by Theorem~\ref{thm:14.3}, $G$ and $H$ have $r$ and $s$ conjugacy classes, respectively, we conclude that $G \times H$ has $rs$ conjugacy classes, with each class of $G \times H$ being the product of a class of $G$ and a class of $H$; hence by Theorem~\ref{thm:14.3}, $G \times H$ has $rs$ irreducible characters.

Let $S_1, \ldots, S_r$ and $T_1, \ldots, T_s$ be the distinct simple $\C G$- and $\C H$-modules, respectively. For each $i$ and $j$, we give $S_i \otimes T_j$ the structure of a $\C(G \times H)$-module by $(g, h)(s \otimes t) = gs \otimes ht$ and linear extension. Let $\tau_{ij}$ be the character of $S_i \otimes T_j$ for each $i$ and $j$. By imitating the proof of part (i) of Proposition~\ref{prop:14.8}, we find that $\tau_{ij}(g, h) = \chi_i(g)\psi_j(h)$; we adopt the notation $\tau_{ij} = \chi_i \times \psi_j$. We note that the $\tau_{ij}$ are distinct and that we can recover the $\chi_i$ and $\psi_j$ from the $\tau_{ij}$ by appropriate restriction. Now for any $i, i', j, j'$ we have
\begin{align*}
    (\tau_{ij}, \tau_{i'j'}) &= \frac{1}{|G \times H|} \sum_{(g,h) \in G \times H} \tau_{ij}(g,h)\overline{\tau_{i'j'}(g,h)} \\
    &= \frac{1}{|G||H|} \sum_{g \in G} \sum_{h \in H} \chi_i(g)\psi_j(h)\overline{\chi_{i'}(g)\psi_{j'}(h)} \\
    &= \left(\frac{1}{|G|} \sum_{g \in G} \chi_i(g)\overline{\chi_{i'}(g)}\right)\left(\frac{1}{|H|} \sum_{h \in H} \psi_j(h)\overline{\psi_{j'}(h)}\right) \\
    &= (\chi_i, \chi_{i'})(\psi_j, \psi_{j'}) = \delta_{ii'}\delta_{jj'}
\end{align*}
by \nameref{thm:row-orthogonality}. In particular, we have $(\tau_{ij}, \tau_{ij}) = 1$ for each $i$ and $j$, which by Corollary~\ref{cor:15.3} implies that $\tau_{ij}$ is an irreducible character. Therefore the $\tau_{ij}$ are the $rs$ irreducible characters of $G \times H$; that is, the set of irreducible characters of $G \times H$ is $\{\chi_i \times \psi_j\}_{i,j}$.
\end{example}

\begin{example}
\label{ex:15.eg.3}
Suppose that $G$ is abelian. Then every conjugacy class of $G$ contains exactly one element, and consequently $G$ has $|G|$ irreducible characters by Theorem~\ref{thm:14.3}. But $\sum_{i=1}^{|G|} f_i^2 = |G|$ by Corollary~\ref{cor:14.2}, so we must have $f_i = 1$ for each $i$. Therefore all irreducible characters of $G$ are linear. By Corollary~\ref{cor:8.8}, $G$ is a direct product of cyclic $p$-groups; let these cyclic $p$-groups be of orders $p_1^{a_1}, \ldots, p_t^{a_t}$ with respective generators $g_1, \ldots, g_t$. We could determine the character table of $G$ from those of the cyclic $p$-groups using Examples~\ref{ex:15.eg.1} and~\ref{ex:15.eg.2}, but there is an alternate method, which we now describe. A linear character $\chi$ of $G$ is just a homomorphism from $G$ to $\C^\times$, and hence to define such a $\chi$ it suffices to specify $\chi(g_i)$ for each $i$, with the only restriction on $\chi(g_i)$ being that it is a $p_i^{a_i}$-th root of unity. Therefore, there is a bijective correspondence between irreducible characters of $G$ and ordered $t$-tuples $(\lambda_1, \ldots, \lambda_t)$ in which each $\lambda_i$ is a $p_i^{a_i}$-th root of unity.

As a simple example, let $G = \langle a, b \rangle \cong \Z_2 \times \Z_2$. Both $a$ and $b$ have order 2, so by the above paragraph we see that the four irreducible characters of $G$ correspond to the ordered pairs $(1,1), (1,-1), (-1,1)$, and $(-1,-1)$. Therefore, the character table for $G$ is
\[
    \begin{array}{c|cccc}
        & 1 & 1 & 1 & 1 \\
        & 1 & a & b & ab \\
        \hline
        \chi_1 & 1 & 1 & 1 & 1 \\
        \chi_2 & 1 & 1 & -1 & -1 \\
        \chi_3 & 1 & -1 & 1 & -1 \\
        \chi_4 & 1 & -1 & -1 & 1
    \end{array}
\]
Here the second and third columns correspond to the ordered pairs given above, and the fourth column is wholly determined by the previous two columns.
\end{example}

\begin{example}
\label{ex:15.eg.4}
Let $X$ be a finite $G$-set. We saw in Section~12 that the vector space $\C X$ having basis $X$ is a $\C G$-module. We wish to determine the character $\pi$ of $\C X$. Consider the matrix, with respect to the basis $X$, of the linear transformation of $\C X$ defined by $g \in G$. Since $gx \in X$ for every $x \in X$, this matrix is a permutation matrix; but $\pi(g)$ is the sum of the diagonal entries of this matrix, so we conclude that $\pi(g)$ equals the number of elements of $X$ which are fixed by $g$.

Let $\chi_1$ be the principal character of $G$. We have
\begin{align*}
    (\pi, \chi_1) &= \frac{1}{|G|}\sum_{g \in G} \pi(g)\chi_1(g) = \frac{1}{|G|}\sum_{g \in G} \pi(g) \\
    &= \frac{1}{|G|}\sum_{g \in G} |\{x \in X \mid gx = x\}| \\
    &= \frac{1}{|G|}|\{(g, x) \in G \times X \mid gx = x\}| \\
    &= \frac{1}{|G|}\sum_{x \in X} |\{g \in G \mid gx = x\}|.
\end{align*}
For each $x \in X$, let $G_x$ be the stabilizer of $x$; using Corollary~\ref{cor:3.5}, we now have
\[
    (\pi, \chi_1) = \sum_{x \in X} \frac{|G_x|}{|G|} = \sum_{\text{orbits } \mathcal{O}} \sum_{x \in \mathcal{O}} \frac{1}{|G : G_x|} = \sum_{\text{orbits } \mathcal{O}} \sum_{x \in \mathcal{O}} \frac{1}{|\mathcal{O}|} = \sum_{\text{orbits } \mathcal{O}} |\mathcal{O}| \cdot \frac{1}{|\mathcal{O}|} = \sum_{\text{orbits } \mathcal{O}} 1,
\]
which allows us to conclude that $(\pi, \chi_1)$ is equal to the number of orbits of $X$ under the action of $G$. (This result is often attributed to Burnside, although as argued in [21] it would be more accurate to call this result the Cauchy-Frobenius lemma.) In particular, if $G$ acts transitively on $X$, then we see via Corollary~\ref{cor:15.4} that the unique expression of the character $\pi - \chi_1$ as a linear combination of the irreducible characters of $G$ does not involve $\chi_1$.
\end{example}

\begin{example}
\label{ex:15.eg.5}
The group $G/G'$ is abelian by Proposition~\ref{prop:2.6}, and hence (as seen in Example~\ref{ex:15.eg.3}) all of its irreducible characters are linear. It follows from Lemma~\ref{lem:15.6} that each linear character of $G/G'$ can be lifted to a linear character of $G$, and that the lifts of distinct characters are distinct. Now let $\chi$ be a linear character of $G$, so that $\chi\colon G \to \C^\times$ is a homomorphism whose kernel is $K_\chi$. Then $G/K_\chi$ is abelian, being isomorphic with $X$, a subgroup of $\C^\times$, and hence $G' \leq K_\chi$ by Proposition~\ref{prop:2.6}. We can now define a linear character $\psi$ of $G/G'$ whose lift is $\chi$ by $\psi(gG') = \chi(g)$. We conclude that every linear character of $G$ is the lift of a linear character of $G/G'$.
\end{example}

\begin{example}
\label{ex:15.eg.6}
Consider $S_3$. The conjugacy classes of $S_3$ are
\[
    \{1\}, \qquad
    \{(1\,2), (1\,3), (2\,3)\}, \qquad
    \{(1\,2\,3), (1\,3\,2)\}.
\]
This follows from Proposition~\ref{prop:1.10}. We have $S_3/A_3 \cong \Z_2$, so $S_3$ has two linear characters, which are lifted from the characters of $\Z_2$ as in Lemma~\ref{lem:15.6}. Letting these characters be $\chi_1$ and $\chi_2$, we have
\[
    \begin{array}{c|ccc}
        & 1 & 3 & 2 \\
        & 1 & (1\,2) & (1\,2\,3) \\
        \hline
        \chi_1 & 1 & 1 & 1 \\
        \chi_2 & 1 & -1 & 1 \\
        \chi_3 & & &
    \end{array}
\]
Since $6 = |S_3| = f_1^2 + f_2^2 + f_3^2 = 2 + f_3^2$ by Corollary~\ref{cor:14.2}, we have $f_3 = 2$. \nameref{thm:col-orthogonality} now gives
\[
    0 = \sum_{i=1}^{3} f_i \chi_i((1\,2)) = 1 \cdot 1 + 1 \cdot (-1) + 2\chi_3((1\,2))
\]
and
\[
    0 = \sum_{i=1}^{3} f_i \chi_i((1\,2\,3)) = 1 \cdot 1 + 1 \cdot 1 + 2\chi_3((1\,2\,3)),
\]
from which we see that the character table of $S_3$ is
\[
    \begin{array}{c|ccc}
        & 1 & 3 & 2 \\
        & 1 & (1\,2) & (1\,2\,3) \\
        \hline
        \chi_1 & 1 & 1 & 1 \\
        \chi_2 & 1 & -1 & 1 \\
        \chi_3 & 2 & 0 & -1
    \end{array}
\]
\end{example}

\begin{example}
\label{ex:15.eg.7}
Consider $S_4$. Recall that $S_4$ has a normal subgroup of order 4, namely
\[
    K = \{1,\, (1\,2)(3\,4),\, (1\,3)(2\,4),\, (1\,4)(2\,3)\}.
\]
The subgroup $S_3$ of $S_4$ intersects $K$ trivially, and $|S_4| = |S_3||K|$; therefore $S_4 = K \rtimes S_3$. In particular, we have $S_4/K \cong S_3$, and thus we can obtain irreducible characters of $S_4$ from the irreducible characters of $S_3$ as in Lemma~\ref{lem:15.6}.

By Proposition~\ref{prop:1.10}, each conjugacy class of $S_4$ consists of all elements having a given cycle structure. Hence $S_4$ has 5 classes, with representatives $1$, $(1\,2)(3\,4)$, $(1\,2\,3)$, $(1\,2)$, and $(1\,2\,3\,4)$; the orders of the classes are 1, 3, 8, 6, and 6, respectively.

Let $\chi_1, \chi_2$, and $\chi_3$ be the irreducible characters of $S_4$ obtained from those of $S_3$. To determine these characters, we need information about the images in $S_4/K$ of the class representatives. We easily see that the images of 1 and $(1\,2)(3\,4)$ are trivial, that the image of $(1\,2\,3)$ has order 3, and that the images of $(1\,2)$ and $(1\,2\,3\,4)$ have order 2.

Using this observation and the character table for $S_3$ obtained in Example~\ref{ex:15.eg.6}, we obtain the following partial character table for $S_4$:
\[
    \begin{array}{c|ccccc}
        & 1 & 3 & 8 & 6 & 6 \\
        & 1 & (1\,2)(3\,4) & (1\,2\,3) & (1\,2) & (1\,2\,3\,4) \\
        \hline
        \chi_1 & 1 & 1 & 1 & 1 & 1 \\
        \chi_2 & 1 & 1 & 1 & -1 & -1 \\
        \chi_3 & 2 & 2 & -1 & 0 & 0 \\
        \chi_4 & & & & & \\
        \chi_5 & & & & &
    \end{array}
\]

Now $S_4$ acts transitively on the set $X = \{1, 2, 3, 4\}$; let $\pi$ be the character of the $\C G$-module $\C X$. Since $\pi(g)$ equals the number of fixed points under the action of $g$ on $X$ by Example~\ref{ex:15.eg.4}, we have
\[
    \begin{array}{c|ccccc}
        & 1 & 3 & 8 & 6 & 6 \\
        & 1 & (1\,2)(3\,4) & (1\,2\,3) & (1\,2) & (1\,2\,3\,4) \\
        \hline
        \pi & 4 & 0 & 1 & 2 & 0 \\
        \pi - \chi_1 & 3 & -1 & 0 & 1 & -1
    \end{array}
\]
We also have
\[
    (\pi - \chi_1, \pi - \chi_1) = \frac{1}{|S_4|} \sum_{i=1}^{5} k_i [(\pi - \chi_1)(g_i)]^2 = \frac{1}{24}(1 \cdot 9 + 3 \cdot 1 + 8 \cdot 0 + 6 \cdot 1 + 6 \cdot 1) = 1,
\]
and hence $\pi - \chi_1$ is irreducible by Corollary~\ref{cor:15.3}. Let $\chi_4 = \pi - \chi_1$. Now $\chi_2 \chi_4 \neq \chi_4$, and $\chi_2 \chi_4$ is irreducible by Proposition~\ref{prop:15.5}; therefore $\chi_5 = \chi_2 \chi_4$, and hence the character table of $S_4$ is
\[
    \begin{array}{c|ccccc}
        & 1 & 3 & 8 & 6 & 6 \\
        & 1 & (1\,2)(3\,4) & (1\,2\,3) & (1\,2) & (1\,2\,3\,4) \\
        \hline
        \chi_1 & 1 & 1 & 1 & 1 & 1 \\
        \chi_2 & 1 & 1 & 1 & -1 & -1 \\
        \chi_3 & 2 & 2 & -1 & 0 & 0 \\
        \chi_4 & 3 & -1 & 0 & 1 & -1 \\
        \chi_5 & 3 & -1 & 0 & -1 & 1
    \end{array}
\]
\end{example}

\begin{example}
\label{ex:15.eg.8}
Let $G$ be a non-abelian group of order 8. We will show that this information completely determines the character table of $G$. Since there are two isomorphism classes of non-abelian groups of order 8 (see page~79), this will imply that the character table does not specify a group up to isomorphism.

We have $|Z(G)| = 2$ by Theorem~\ref{thm:8.1} and Lemma~\ref{lem:8.2}, and hence $Z(G) \cong \Z_2$. Since $G$ is non-abelian and $G/Z(G)$ is abelian (being a group of order 4), this forces $G' = Z(G)$ by Proposition~\ref{prop:2.6}. Now $|G/G'| = 4$, and by Lemma~\ref{lem:8.2} $G/G'$ cannot be cyclic since $G$ is non-abelian and $G' = Z(G)$, so we must have $G/G' \cong \Z_2 \times \Z_2$. By Examples~\ref{ex:15.eg.3} and~\ref{ex:15.eg.5}, we now know that $G$ has exactly 4 linear characters. By Corollary~\ref{cor:14.2}, we have $\sum_{i=1}^{r} f_i^2 = 8$, where $f_i = 1$ for $1 \leq i \leq 4$ and $f_i > 1$ for $i > 4$; this forces $r = 5$ and $f_5 = 2$.

Let $x$ be the generator of $G'$. As $G' = Z(G)$, we see that 1 and $x$ are the only elements of $G$ lying in single-element conjugacy classes, and hence each of the other three classes must have order 2. Let $a$, $b$, and $c$ be representatives of the multi-element classes. Since $G/G'$ also has four conjugacy classes, we see that the images of $a$, $b$, and $c$ in $G/G'$ must lie in distinct conjugacy classes. Using the character table for $\Z_2 \times \Z_2$ determined in Example~\ref{ex:15.eg.3}, we can now obtain a partial character table for $G$ by lifting:
\[
    \begin{array}{c|ccccc}
        & 1 & 1 & 2 & 2 & 2 \\
        & 1 & x & a & b & c \\
        \hline
        \chi_1 & 1 & 1 & 1 & 1 & 1 \\
        \chi_2 & 1 & 1 & 1 & -1 & -1 \\
        \chi_3 & 1 & 1 & -1 & 1 & -1 \\
        \chi_4 & 1 & 1 & -1 & -1 & 1 \\
        \chi_5 & 2 & & & &
    \end{array}
\]
We see via \nameref{thm:col-orthogonality} that $\chi_5(a) = \chi_5(b) = \chi_5(c) = 0$ and that $\chi_5(x) = -2$, completing the character table.
\end{example}

\begin{example}
\label{ex:15.eg.9}
Consider the group $A_5$. Using Proposition~\ref{prop:1.10}, we find that the elements of $A_5$ form 4 conjugacy classes in $S_5$, with representatives $1$, $(1\,2)(3\,4)$, $(1\,2\,3)$, and $(1\,2\,3\,4\,5)$ and respective orders 1, 15, 20, and 24. However, elements of $A_5$ that are conjugate in $S_5$ may not be conjugate in $A_5$.

Let $x \in A_5$, and let $\mathcal{K}$ be the conjugacy class in $S_5$ of $x$. We see from Proposition~\ref{prop:3.13} that $|\mathcal{K}| = |S_5 : C_{S_5}(x)|$ and that the conjugacy class in $A_5$ of $x$ is contained in $\mathcal{K}$ and has order $|A_5 : C_{A_5}(x)|$. Suppose first that $C_{S_5}(x) \not\leq A_5$. This forces $A_5 C_{S_5}(x) = S_5$ since $A_5$ is a maximal subgroup of $S_5$. The first isomorphism theorem gives $S_5/A_5 = A_5 C_{S_5}(x)/A_5 \cong C_{S_5}(x)/C_{S_5}(x) \cap A_5 = C_{S_5}(x)/C_{A_5}(x)$, and hence
\begin{align*}
    |A_5 : C_{A_5}(x)| &= \frac{|S_5 : C_{A_5}(x)|}{|S_5 : A_5|} = \frac{|S_5 : C_{A_5}(x)|}{|C_{S_5}(x) : C_{A_5}(x)|} = |S_5 : C_{S_5}(x)| \\
    &= |\mathcal{K}|.
\end{align*}
Thus $\mathcal{K}$ is the conjugacy class in $A_5$ of $x$. Now suppose $C_{S_5}(x) \leq A_5$. Then we have $C_{A_5}(x) = C_{S_5}(x) \cap A_5 = C_{S_5}(x)$, and hence
\[
    |A_5 : C_{A_5}(x)| = |A_5 : C_{S_5}(x)| = \tfrac{1}{2}|S_5 : C_{S_5}(x)| = \tfrac{1}{2}|\mathcal{K}|,
\]
from which we see that $\mathcal{K}$ splits into two conjugacy classes in $A_5$ of equal cardinality.

We easily see that each of the elements 1, $(1\,2)(3\,4)$, and $(1\,2\,3)$ commutes with some odd permutation in $S_5$; for example, $(4\,5)$ and $(1\,2\,3)$ commute. Hence by the above paragraph, the conjugacy classes in $A_5$ of these elements are the respective classes in $S_5$. Since $24 \nmid 60$, we can now conclude that the conjugacy class in $S_5$ of $(1\,2\,3\,4\,5)$ splits into two classes in $A_5$, and hence that the conjugacy classes of $A_5$ have orders 1, 15, 20, 12, and 12. (Using this fact and the remark made after Proposition~\ref{prop:3.13}, we could now give a slick proof of the simplicity of $A_5$.) We will now show that $x = (1\,2\,3\,4\,5)$ and $x^2 = (1\,3\,5\,2\,4)$ are not conjugate in $A_5$, thereby establishing that the conjugacy class in $A_5$ of $x$ splits in $A_5$ into the class containing $x$ and the class containing $x^2$.

Suppose that $g \in A_5$ is such that $gxg^{-1} = x^2$. Then we have $g\langle x \rangle g^{-1} = \langle x^2 \rangle = \langle x \rangle$, and hence $g \in N_{A_5}(\langle x \rangle)$. Now $\langle x \rangle$ is a Sylow 5-subgroup of $A_5$, and we observed in the proof of Theorem~\ref{thm:7.4} that $A_5$ has 6 Sylow 5-subgroups. Thus $|N_{A_5}(\langle x \rangle)| = 10$, and we find that $N_{A_5}(\langle x \rangle) = \langle x \rangle \rtimes \langle t \rangle \cong D_{10}$, where $t = (2\,5)(3\,4)$. But $txt^{-1} = x^{-1} \neq x^2$; from this, we see that there is no element $g \in N_{A_5}(\langle x \rangle)$ such that $gxg^{-1} = x^2$. Contradiction; therefore $x$ and $x^2$ are not conjugate in $A_5$.

We now see that $A_5$ has 5 conjugacy classes, with representatives $1$, $(1\,2)(3\,4)$, $(1\,2\,3)$, $(1\,2\,3\,4\,5)$, and $(1\,3\,5\,2\,4)$ and respective orders 1, 15, 20, 12, and 12. Since $A_5$ is simple, we cannot produce irreducible characters of $A_5$ by lifting irreducible characters of quotient groups, as we have done in previous examples. However, $A_5$ is a group of permutations, so we have as a starting point the characters of certain obvious permutation modules.

Let $X = \{1, 2, 3, 4, 5\}$; $A_5$ acts transitively on $X$ in the obvious way. Let $\pi$ be the character of the $\C A_5$-module $\C X$. By Example~\ref{ex:15.eg.4}, we see that $\pi(g)$ equals the number of fixed points of $X$ under the action of $g \in A_5$, and hence we have
\[
    \begin{array}{c|ccccc}
        & 1 & 15 & 20 & 12 & 12 \\
        & 1 & (1\,2)(3\,4) & (1\,2\,3) & (1\,2\,3\,4\,5) & (1\,3\,5\,2\,4) \\
        \hline
        \pi & 5 & 1 & 2 & 0 & 0 \\
        \pi - \chi_1 & 4 & 0 & 1 & -1 & -1
    \end{array}
\]
We also have
\[
    (\pi - \chi_1, \pi - \chi_1) = \frac{1}{|A_5|}\sum_{i=1}^{5} k_i[(\pi - \chi_1)(g_i)]^2 = \frac{1}{60}(1 \cdot 16 + 15 \cdot 0 + 20 \cdot 1 + 12 \cdot 1 + 12 \cdot 1) = 1,
\]
and hence $\pi - \chi_1$ is irreducible by Corollary~\ref{cor:15.3}. Let $\chi_2 = \pi - \chi_1$.

Let $Y$ be the set of two-element subsets of $X$; this is a transitive $A_5$-set under the obvious action. Let $\psi$ be the character of the $\C A_5$-module $\C Y$. We have
\[
    \begin{array}{c|ccccc}
        & 1 & 15 & 20 & 12 & 12 \\
        & 1 & (1\,2)(3\,4) & (1\,2\,3) & (1\,2\,3\,4\,5) & (1\,3\,5\,2\,4) \\
        \hline
        \psi & 10 & 2 & 1 & 0 & 0 \\
        \psi - \chi_1 & 9 & 1 & 0 & -1 & -1
    \end{array}
\]
Now
\[
    (\psi - \chi_1, \psi - \chi_1) = \frac{1}{|A_5|}\sum_{i=1}^{5} k_i[(\psi - \chi_1)(g_i)]^2 = \frac{1}{60}(1 \cdot 81 + 15 \cdot 1 + 20 \cdot 0 + 12 \cdot 1 + 12 \cdot 1) = 2,
\]
so by Corollary~\ref{cor:15.3}, $\psi - \chi_1$ is the sum of two irreducible characters; as noted in Example~\ref{ex:15.eg.4}, neither of these two characters is $\chi_1$, since $Y$ is transitive. We have
\[
    (\psi - \chi_1, \chi_2) = \frac{1}{|A_5|}\sum_{i=1}^{5} k_i(\psi - \chi_1)(g)\chi_2(g) = \frac{1}{60}(1 \cdot 9 \cdot 4 + 15 \cdot 1 \cdot 0 + 20 \cdot 0 \cdot 1 + 24 \cdot (-1)(-1)) = 1,
\]
and therefore by Corollary~\ref{cor:15.4} we see that $\psi - \chi_1 - \chi_2$ is an irreducible character that is neither $\chi_1$ nor $\chi_2$. Let $\chi_3 = \psi - \chi_1 - \chi_2$, and observe that $f_3 = 5$.

By Corollary~\ref{cor:14.2} we have $\sum_{i=1}^{5} f_i^2 = |A_5| = 60$, which gives $f_4^2 + f_5^2 = 18$ and hence $f_4 = f_5 = 3$. Thus we have the following partial character table:
\[
    \begin{array}{c|ccccc}
        & 1 & 15 & 20 & 12 & 12 \\
        & 1 & (1\,2)(3\,4) & (1\,2\,3) & (1\,2\,3\,4\,5) & (1\,3\,5\,2\,4) \\
        \hline
        \chi_1 & 1 & 1 & 1 & 1 & 1 \\
        \chi_2 & 4 & 0 & 1 & -1 & -1 \\
        \chi_3 & 5 & 1 & -1 & 0 & 0 \\
        \chi_4 & 3 & & & & \\
        \chi_5 & 3 & & & &
    \end{array}
\]

Let $a = \chi_4((1\,2)(3\,4))$ and $b = \chi_5((1\,2)(3\,4))$. By \nameref{thm:col-orthogonality}, we have
\[
    0 = \sum_{i=1}^{5} f_i \chi_i((1\,2)(3\,4)) = 1 \cdot 1 + 4 \cdot 0 + 5 \cdot 1 + 3a + 3b,
\]
and hence $a + b = -2$. By part (iii) of Proposition~\ref{prop:14.4}, we see that each of $a$ and $b$ is a sum of three square roots of unity; therefore $a, b \in \{-3, -1, 1, 3\}$. But by \nameref{thm:col-orthogonality} we have
\[
    4 = \frac{60}{15} = \frac{|A_5|}{k_2} = \sum_{i=1}^{5} |\chi_i((1\,2)(3\,4))|^2 = 1 + 0 + 1 + a^2 + b^2,
\]
and hence $a^2 + b^2 = 2$. It now follows that $a = b = -1$.

By \nameref{thm:col-orthogonality}, we have
\[
    3 = \frac{60}{20} = \frac{|A_5|}{k_3} = \sum_{i=1}^{5} |\chi_i((1\,2\,3))|^2 = 1 + 1 + 1 + |\chi_4((1\,2\,3))|^2 + |\chi_5((1\,2\,3))|^2;
\]
thus $|\chi_4((1\,2\,3))|^2 + |\chi_5((1\,2\,3))|^2 = 0$, and therefore we conclude that $\chi_4((1\,2\,3)) = \chi_5((1\,2\,3)) = 0$. We now have the following partial character table:
\[
    \begin{array}{c|ccccc}
        & 1 & 15 & 20 & 12 & 12 \\
        & 1 & (1\,2)(3\,4) & (1\,2\,3) & (1\,2\,3\,4\,5) & (1\,3\,5\,2\,4) \\
        \hline
        \chi_1 & 1 & 1 & 1 & 1 & 1 \\
        \chi_2 & 4 & 0 & 1 & -1 & -1 \\
        \chi_3 & 5 & 1 & -1 & 0 & 0 \\
        \chi_4 & 3 & -1 & 0 & & \\
        \chi_5 & 3 & -1 & 0 & &
    \end{array}
\]

Let $x = (1\,2\,3\,4\,5)$. Recall that in proving that $x$ and $x^2$ are not conjugate in $A_5$, we found that $x$ and $x^{-1}$ are conjugate in $A_5$; hence $\chi_4(x) = \chi_4(x^{-1})$, and thus $\chi_4(x)$ is real-valued by part (iv) of Proposition~\ref{prop:14.4}. The same argument shows that $\chi_4(x^2)$, $\chi_5(x)$, and $\chi_5(x^2)$ are also real-valued. Let $c = \chi_4(x)$ and $d = \chi_4(x^2)$. Then by \nameref{thm:row-orthogonality}, we have
\[
    0 = \sum_{i=1}^{5} k_i \chi_4(g_i) = 1 \cdot 3 + 15 \cdot (-1) + 20 \cdot 0 + 12c + 12d,
\]
which gives $c + d = 1$, and also (since $\chi_4$ is real-valued)
\[
    60 = \sum_{i=1}^{5} k_i \chi_4(g_i)^2 = 1 \cdot 9 + 15 \cdot 1 + 20 \cdot 0 + 12c^2 + 12d^2,
\]
which gives $c^2 + d^2 = 3$. Hence $c^2 - c - 1 = 0$, so without loss of generality we have $c = \frac{1 + \sqrt{5}}{2}$ and $d = 1 - c = \frac{1 - \sqrt{5}}{2}$. We can similarly show that we must have $\chi_5(x), \chi_5(x^2) \in \{\frac{1 + \sqrt{5}}{2}, \frac{1 - \sqrt{5}}{2}\}$; we leave it to the reader to verify that $\chi_5(x) = \frac{1 - \sqrt{5}}{2}$ and $\chi_5(x^2) = \frac{1 + \sqrt{5}}{2}$, which completes the character table of $A_5$.
\end{example}

\subsection*{Exercises}

Throughout these exercises, $G$ denotes a finite group, and $\mathcal{X}$ denotes the character table of $G$.

\begin{exercise}
\label{ex:15.1}
What is the determinant of $\mathcal{X}$?
\end{exercise}
\begin{solution}
    Let $g_1,\ldots,g_r$ be the chosen conjugacy class representatives. By the \nameref{thm:col-orthogonality},
    \[
        \mathcal{X}^{T}\overline{\mathcal{X}}
        =
        \operatorname{diag}\bigl(|C_G(g_1)|,\ldots,|C_G(g_r)|\bigr).
    \]
    Taking determinants gives
    \[
        |\det(\mathcal{X})|^2
        =
        \det(\mathcal{X})\overline{\det(\mathcal{X})}
        =
        \prod_{j=1}^{r}|C_G(g_j)|.
    \]
    Hence
    \[
        |\det(\mathcal{X})|
        =
        \left(\prod_{j=1}^{r}|C_G(g_j)|\right)^{1/2}
        =
        \frac{|G|^{r/2}}{(k_1\cdots k_r)^{1/2}}.
    \]
    We can also determine whether the determinant is real or purely imaginary. Let $m$ be the number of unordered pairs of non-real complex-conjugate irreducible characters. By Exercise~\ref{ex:14.6}, complex conjugation permutes the rows of $\mathcal{X}$, fixing every real-valued row and interchanging the two rows in each of these $m$ pairs. Therefore
    \[
        \overline{\det(\mathcal{X})}=(-1)^m\det(\mathcal{X}),
    \]
    and so
    \[
        \det(\mathcal{X})^2=(-1)^m\prod_{j=1}^{r}|C_G(g_j)|.
    \]
    Thus, up to the unavoidable sign caused by reordering rows or columns,
    \[
        \det(\mathcal{X})
        =
        \pm i^m\left(\prod_{j=1}^{r}|C_G(g_j)|\right)^{1/2}.
    \]
\end{solution}

\begin{exercise}
\label{ex:15.2}
Show that the sum of the entries in any row of $\mathcal{X}$ is a non-negative integer.
\end{exercise}
\begin{solution}
    Let $G$ act on itself by conjugation, and let $\pi$ be the character of the corresponding permutation module. By Example~\ref{ex:15.eg.4}, $\pi(g)$ is the number of elements of $G$ fixed by $g$. Since
    \[
        gxg^{-1}=x
        \quad\Longleftrightarrow\quad
        gx=xg,
    \]
    these fixed points are exactly the elements of $C_G(g)$, and hence $\pi(g)=|C_G(g)|$. Fix an irreducible character $\chi_i$, and let $s_i=\sum_{j=1}^{r}\chi_i(g_j)$ be the sum of the entries in its row. Grouping the inner product $(\pi,\chi_i)$ by conjugacy classes gives
    \begin{align*}
        (\pi,\chi_i)
        &=\frac{1}{|G|}\sum_{j=1}^{r}k_j|C_G(g_j)|\overline{\chi_i(g_j)} \\
        &=\sum_{j=1}^{r}\overline{\chi_i(g_j)}
        =\overline{s_i},
    \end{align*}
    because $k_j|C_G(g_j)|=|G|$. By Corollary~\ref{cor:15.4}, $(\pi,\chi_i)$ is the multiplicity of $\chi_i$ in $\pi$, so it is a non-negative integer. It is therefore real, and
    \[
        s_i=\overline{(\pi,\chi_i)}=(\pi,\chi_i)\in\Z_{\geq0}.
    \]
    Hence the sum of the entries in every row of $\mathcal{X}$ is a non-negative integer.
\end{solution}

\begin{exercise}
\label{ex:15.3}
Compute the character of the $G$-set $G$, where $G$ acts via conjugation, and determine the multiplicities in this character of the irreducible characters.
\end{exercise}
\begin{solution}
    Let $\pi$ be the character of the permutation module corresponding to the conjugation action of $G$ on itself. As in Exercise~\ref{ex:15.2}, an element $x\in G$ is fixed by $g$ exactly when $x\in C_G(g)$. Therefore
    \[
        \pi(g)=|C_G(g)|,
        \qquad
        \pi(g_j)=|C_G(g_j)|=\frac{|G|}{k_j}.
    \]
    Let $m_i$ be the multiplicity of $\chi_i$ in $\pi$. By Corollary~\ref{cor:15.4},
    \begin{align*}
        m_i
        &=(\pi,\chi_i) \\
        &=\frac{1}{|G|}\sum_{j=1}^{r}k_j\pi(g_j)\overline{\chi_i(g_j)} \\
        &=\sum_{j=1}^{r}\overline{\chi_i(g_j)}.
    \end{align*}
    By Exercise~\ref{ex:15.2}, the sum of the entries in the $\chi_i$-row is a non-negative integer and is therefore real. Hence
    \[
        m_i=\sum_{j=1}^{r}\chi_i(g_j).
    \]
    We conclude that the conjugation character is
    \[
        \pi
        =
        \sum_{i=1}^{r}\left(\sum_{j=1}^{r}\chi_i(g_j)\right)\chi_i.
    \]
\end{solution}

\begin{exercise}
\label{ex:15.4}
Observe that the complex conjugate of an irreducible character is an irreducible character. (This is either Exercise~\ref{ex:14.6} or an easy consequence of part (ii) of Proposition~\ref{prop:14.8} and Corollary~\ref{cor:15.3}.) Let $P$ be the $r \times r$ permutation matrix such that $P\mathcal{X}$ is the matrix obtained from $\mathcal{X}$ by interchanging two rows if their corresponding characters are conjugate. Let $Q$ be the $r \times r$ permutation matrix such that $\mathcal{X}Q$ is the matrix obtained from $\mathcal{X}$ by interchanging two columns if their corresponding conjugacy classes are inverse, in the sense that one class consists of the inverses of the elements of the other class. Show that $P\mathcal{X} = \mathcal{X}Q$.
\end{exercise}
\begin{solution}
    Let $p$ be the permutation of $\{1,\ldots,r\}$ determined by $\chi_{p(i)}=\overline{\chi_i}$, and let $q$ be the permutation determined by requiring that $g_{q(j)}$ be conjugate to $g_j^{-1}$. By the definitions of $P$ and $Q$, we have
    \[
        (P\mathcal{X})_{ij}
        =
        \chi_{p(i)}(g_j)
        =
        \overline{\chi_i(g_j)}
    \]
    and
    \[
        (\mathcal{X}Q)_{ij}
        =
        \chi_i(g_{q(j)})
        =
        \chi_i(g_j^{-1}).
    \]
    The last equality holds because characters are constant on conjugacy classes. By part (iv) of Proposition~\ref{prop:14.4},
    \[
        \chi_i(g_j^{-1})=\overline{\chi_i(g_j)}.
    \]
    Hence $(P\mathcal{X})_{ij}=(\mathcal{X}Q)_{ij}$ for every $i$ and $j$, and therefore $P\mathcal{X}=\mathcal{X}Q$.
\end{solution}

\begin{exercise}
\label{ex:15.5}
Suppose that $|G|$ is odd. Show that the only irreducible character that is real-valued is the principal character.
\end{exercise}
\begin{solution}
    By the \nameref{thm:row-orthogonality}, the rows of $\mathcal{X}$ are linearly independent, so $\mathcal{X}$ is invertible. Exercise~\ref{ex:15.4} gives
    \[
        P=\mathcal{X}Q\mathcal{X}^{-1}.
    \]
    Thus $P$ and $Q$ are similar and have the same trace. The trace of a permutation matrix is the number of points fixed by the corresponding permutation. Therefore $\operatorname{tr}(P)$ is the number of real-valued irreducible characters, while $\operatorname{tr}(Q)$ is the number of conjugacy classes that are equal to their inverse classes.\\
    We show that the identity class is the only conjugacy class fixed by inversion. Let $\mathcal{K}$ be such a class. By Proposition~\ref{prop:3.13}, $|\mathcal{K}|$ divides $|G|$, so $|\mathcal{K}|$ is odd. Inversion restricts to an involution on $\mathcal{K}$. Since all non-fixed points of an involution occur in pairs, this involution on the odd-sized set $\mathcal{K}$ has a fixed point $x$. Then $x=x^{-1}$, so $x^2=1$. Since $|G|$ is odd, \nameref{thm:lagrange} forces $x=1$, and therefore $\mathcal{K}=\{1\}$. Consequently $\operatorname{tr}(Q)=1$, so also $\operatorname{tr}(P)=1$.\\
    The principal character is real-valued, and $\operatorname{tr}(P)=1$ shows that it is the only real-valued irreducible character.
\end{solution}

\begin{exercise}
\label{ex:15.6}
Let $\chi$ be an irreducible character of $G$, and let $e \in \C G$ be the identity element of the corresponding matrix summand of $\C G$. Show that
\[
    e = \frac{\chi(1)}{|G|} \sum_{g \in G} \chi(g^{-1}) g.
\]
\end{exercise}
\begin{solution}
    Write
    \[
        e=\sum_{h\in G}a_hh.
    \]
    Fix $g\in G$, and let $L_{eg^{-1}}$ denote left multiplication by $eg^{-1}$ on the regular module $\C G$. For each basis element $x\in G$, we have
    \[
        (eg^{-1})x=\sum_{h\in G}a_hhg^{-1}x.
    \]
    The coefficient of $x$ in this expression is $a_g$, since $hg^{-1}x=x$ iff $h=g$. Thus every diagonal entry of $L_{eg^{-1}}$ is $a_g$, and
    \[
        \operatorname{tr}(L_{eg^{-1}})=|G|a_g.
    \]
    Let $S$ be a simple $\C G$-module affording $\chi$, and put $f=\chi(1)=\dim_{\C}S$. By Theorem~\ref{thm:14.1}, the regular module contains $f$ copies of $S$. Since $e$ is the identity element of the matrix summand corresponding to $S$, it acts as the identity on $S$ and as zero on every non-isomorphic simple module. Consequently, $eg^{-1}$ acts as $g^{-1}$ on each of the $f$ copies of $S$ and as zero on every other simple summand of the regular module. Taking traces gives
    \[
        \operatorname{tr}(L_{eg^{-1}})=f\chi(g^{-1})=\chi(1)\chi(g^{-1}).
    \]
    Comparing the two expressions for $\operatorname{tr}(L_{eg^{-1}})$ yields
    \[
        a_g=\frac{\chi(1)}{|G|}\chi(g^{-1})
    \]
    for every $g\in G$. Substituting these coefficients into the expression for $e$, we obtain
    \[
        e=\frac{\chi(1)}{|G|}\sum_{g\in G}\chi(g^{-1})g.
    \]
\end{solution}

\begin{exercise}
\label{ex:15.7}
If $N \trianglelefteq G$ and $g \in G$, show that $|C_G(g)| \geq |C_{G/N}(gN)|$.
\end{exercise}
\begin{solution}
    Let $\chi_1,\ldots,\chi_r$ be the irreducible characters of $G$. By Lemma~\ref{lem:15.13}, the irreducible characters of $G/N$ correspond exactly to those $\chi_i$ for which $N\leq K_i$, and the corresponding character of $G/N$ takes the value $\chi_i(g)$ at $gN$. Applying the \nameref{thm:col-orthogonality} first in $G/N$ and then in $G$, we obtain
    \[
        |C_{G/N}(gN)|
        =
        \sum_{\substack{1\leq i\leq r\\N\leq K_i}}|\chi_i(g)|^2
        \leq
        \sum_{i=1}^{r}|\chi_i(g)|^2
        =
        |C_G(g)|.
    \]
    The inequality holds because every summand is non-negative. Therefore $|C_G(g)|\geq|C_{G/N}(gN)|$.
\end{solution}

\begin{exercise}
\label{ex:15.8}
Show that if $G$ acts doubly transitively on $X$, then the character of $\C X$ is the sum of the principal character and exactly one other irreducible character.
\end{exercise}
\begin{solution}
    As written, the assertion requires the additional assumption $|X|\geq2$. If $|X|=1$, the action is doubly transitive under the definition in Section~3, but $\C X$ affords only the principal character. We therefore prove the intended statement under the assumption $|X|\geq2$.\\
    Let $\pi$ be the character of $\C X$, and let $G$ act diagonally on $X\times X$. By Example~\ref{ex:15.eg.4}, $\pi(g)$ is the number of points of $X$ fixed by $g$. Hence the number of points of $X\times X$ fixed by $g$ is $\pi(g)^2=|\pi(g)|^2$. Applying Example~\ref{ex:15.eg.4} to $X\times X$, we see that
    \[
        (\pi,\pi)
        =
        \frac{1}{|G|}\sum_{g\in G}|\pi(g)|^2
    \]
    is the number of orbits of $G$ on $X\times X$. Double transitivity gives exactly two such orbits: the diagonal
    \[
        \{(x,x)\mid x\in X\}
    \]
    and the set of ordered pairs $(x,y)$ with $x\neq y$. Thus $(\pi,\pi)=2$. By Corollary~\ref{cor:15.3}, $\pi$ is the sum of two distinct irreducible characters.\\
    A doubly transitive action is transitive, so Example~\ref{ex:15.eg.4} gives $(\pi,\chi_1)=1$. Hence one of the two irreducible constituents is the principal character, and the other is a single non-principal irreducible character $\chi$. Therefore
    \[
        \pi=\chi_1+\chi.
    \]
\end{solution}

\begin{exercise}
\label{ex:15.9}
Determine the character table of $D_{2n}$ for any $n \in \N$.
\end{exercise}
\begin{solution}
    Write
    \[
        D_{2n}=\langle r,s\mid r^n=s^2=1,\ srs=r^{-1}\rangle,
    \]
    where $r$ is a rotation and $s$ is a reflection. Every element is either $r^j$ or $sr^j$. Conjugation by $s$ sends $r^j$ to $r^{-j}$, while
    \[
        r^a(sr^j)r^{-a}=sr^{j-2a}.
    \]
    We therefore separate the cases in which $n$ is odd and even.\\
    Suppose first that $n$ is odd. The conjugacy classes are
    \[
        \{1\},
        \qquad
        \{r^j,r^{-j}\}\quad\left(1\leq j\leq\frac{n-1}{2}\right),
        \qquad
        \{sr^a\mid0\leq a<n\}.
    \]
    The last set is one class because multiplication by $2$ is invertible modulo $n$. There are two linear characters $\lambda_{\delta}$, where $\delta\in\{1,-1\}$, given by
    \[
        \lambda_{\delta}(r^j)=1,
        \qquad
        \lambda_{\delta}(sr^j)=\delta.
    \]
    Let $\zeta=e^{2\pi i/n}$. For $1\leq k\leq(n-1)/2$, define a representation by
    \[
        r\longmapsto
        \begin{pmatrix}
            \zeta^k&0\\
            0&\zeta^{-k}
        \end{pmatrix},
        \qquad
        s\longmapsto
        \begin{pmatrix}
            0&1\\
            1&0
        \end{pmatrix}.
    \]
    Its character $\chi_k$ satisfies
    \[
        \chi_k(r^j)=\zeta^{kj}+\zeta^{-kj}=2\cos\left(\frac{2\pi kj}{n}\right),
        \qquad
        \chi_k(sr^j)=0.
    \]
    The two eigenvalues of $r$ in this representation are distinct, and $s$ interchanges their eigenspaces, so the representation has no invariant one-dimensional subspace and is irreducible. The characters $\chi_k$ are pairwise distinct, since the values $2\cos(2\pi k/n)$ at $r$ are distinct in the given range. We have constructed
    \[
        2+\frac{n-1}{2}=\frac{n+3}{2}
    \]
    irreducible characters, exactly the number of conjugacy classes. Hence the character table for odd $n$ is
    \[
        \begin{array}{c|ccc}
            &1&r^{\pm j}\ \left(1\leq j\leq\frac{n-1}{2}\right)&sr^a\ (0\leq a<n)\\
            \hline
            \lambda_{1}&1&1&1\\
            \lambda_{-1}&1&1&-1\\
            \chi_k\ \left(1\leq k\leq\frac{n-1}{2}\right)&2&2\cos\left(\frac{2\pi kj}{n}\right)&0
        \end{array}
    \]
    where the class sizes are respectively $1$, $2$, and $n$.\\
    Now suppose that $n$ is even. The conjugacy classes are
    \[
        \{1\},\quad
        \{r^{n/2}\},\quad
        \{r^j,r^{-j}\}\ \left(1\leq j<\frac n2\right),\quad
        \{sr^{2a}\mid0\leq a<\tfrac n2\},\quad
        \{sr^{2a+1}\mid0\leq a<\tfrac n2\}.
    \]
    There are four linear characters $\lambda_{\varepsilon,\delta}$, indexed by $\varepsilon,\delta\in\{1,-1\}$, defined by
    \[
        \lambda_{\varepsilon,\delta}(r^j)=\varepsilon^j,
        \qquad
        \lambda_{\varepsilon,\delta}(sr^j)=\delta\varepsilon^j.
    \]
    For $1\leq k<n/2$, the same two-dimensional matrices as above define irreducible characters $\chi_k$. Thus the character table for even $n$ is
    \[
        \begin{array}{c|ccccc}
            &1&r^{n/2}&r^{\pm j}\ \left(1\leq j<\frac n2\right)&sr^{2a}&sr^{2a+1}\\
            \hline
            \lambda_{\varepsilon,\delta}&1&\varepsilon^{n/2}&\varepsilon^j&\delta&\delta\varepsilon\\
            \chi_k\ \left(1\leq k<\frac n2\right)&2&2(-1)^k&2\cos\left(\frac{2\pi kj}{n}\right)&0&0
        \end{array}
    \]
    where $\varepsilon,\delta\in\{1,-1\}$ and the class sizes are respectively $1$, $1$, $2$, $n/2$, and $n/2$. These give
    \[
        4+\left(\frac n2-1\right)=\frac n2+3
    \]
    pairwise distinct irreducible characters, again exactly the number of conjugacy classes. The formulas also cover $n=1$ and $n=2$, when the indexed families of two-dimensional characters are empty.
\end{solution}

\begin{exercise}
\label{ex:15.10}
Determine the character table of $S_5$.
\end{exercise}
\begin{solution}
    By Proposition~\ref{prop:1.10}, the conjugacy classes of $S_5$ are determined by cycle type. We order the class representatives and their class sizes as follows:
    \[
        \begin{array}{c|ccccccc}
            &1&15&20&24&10&20&30\\
            &1&(1\,2)(3\,4)&(1\,2\,3)&(1\,2\,3\,4\,5)&(1\,2)&(1\,2\,3)(4\,5)&(1\,2\,3\,4)
        \end{array}
    \]
    Let $\chi_1$ be the principal character and let $\chi_2$ be the sign character. Let $\pi$ be the permutation character for the natural action of $S_5$ on $\{1,2,3,4,5\}$. By Example~\ref{ex:15.eg.4}, $\pi(g)$ is the number of fixed points of $g$, so
    \[
        \pi=(5,1,2,0,3,0,1).
    \]
    The action is transitive, so $(\pi,\chi_1)=1$ by Example~\ref{ex:15.eg.4}; hence $\pi-\chi_1$ is a character. A direct calculation gives
    \[
        (\pi-\chi_1,\pi-\chi_1)
        =
        \frac{1}{120}(1\cdot16+20\cdot1+24\cdot1+10\cdot4+20\cdot1)
        =1.
    \]
    Hence $\chi_3=\pi-\chi_1=(4,0,1,-1,2,-1,0)$ is irreducible by Corollary~\ref{cor:15.3}. Proposition~\ref{prop:15.5} then gives the distinct irreducible character
    \[
        \chi_4=\chi_2\chi_3=(4,0,1,-1,-2,1,0).
    \]
    Next let $S_5$ act by conjugation on its set of ten transpositions, and let $\tau$ be the corresponding permutation character. Counting the transpositions that commute with each class representative gives
    \[
        \tau=(10,2,1,0,4,1,0).
    \]
    This action is transitive, so $(\tau,\chi_1)=1$ by Example~\ref{ex:15.eg.4}. Direct calculations also give
    \[
        (\tau,\chi_3)
        =
        \frac{1}{120}(1\cdot10\cdot4+20\cdot1\cdot1+10\cdot4\cdot2+20\cdot1\cdot(-1))
        =1
    \]
    and
    \[
        (\tau,\tau)
        =
        \frac{1}{120}(1\cdot100+15\cdot4+20\cdot1+10\cdot16+20\cdot1)
        =3.
    \]
    Thus $\tau$ contains one copy each of $\chi_1$ and $\chi_3$, and the remaining character
    \[
        \chi_5
        =
        \tau-\chi_1-\chi_3
        =
        (5,1,-1,0,1,1,-1)
    \]
    has $(\chi_5,\chi_5)=1$. Therefore $\chi_5$ is irreducible, and Proposition~\ref{prop:15.5} gives the distinct irreducible character
    \[
        \chi_6=\chi_2\chi_5=(5,1,-1,0,-1,-1,1).
    \]
    There are seven conjugacy classes, so one irreducible character $\chi_7$ remains. Corollary~\ref{cor:14.2} gives
    \[
        \chi_7(1)^2
        =
        120-(1^2+1^2+4^2+4^2+5^2+5^2)
        =36,
    \]
    and hence $\chi_7(1)=6$. For each non-identity class representative $g$, orthogonality with the identity column gives
    \[
        \sum_{i=1}^{7}\chi_i(1)\chi_i(g)=0.
    \]
    Substituting the six known rows into these equations determines the remaining values, giving
    \[
        \chi_7=(6,-2,0,1,0,0,0).
    \]
    Therefore the character table of $S_5$ is
    \[
        \begin{array}{c|ccccccc}
            &1&15&20&24&10&20&30\\
            &1&(1\,2)(3\,4)&(1\,2\,3)&(1\,2\,3\,4\,5)&(1\,2)&(1\,2\,3)(4\,5)&(1\,2\,3\,4)\\
            \hline
            \chi_1&1&1&1&1&1&1&1\\
            \chi_2&1&1&1&1&-1&-1&-1\\
            \chi_3&4&0&1&-1&2&-1&0\\
            \chi_4&4&0&1&-1&-2&1&0\\
            \chi_5&5&1&-1&0&1&1&-1\\
            \chi_6&5&1&-1&0&-1&-1&1\\
            \chi_7&6&-2&0&1&0&0&0
        \end{array}
    \]
\end{solution}

\end{document}
